\section{Economía de los centros de datos: de los MW de la grid a los useful tokens}

Tener chips no equivale a contar con capacidad de cómputo utilizable. La energía que entra de la red eléctrica pasa por el suministro y la distribución eléctrica, el UPS, el cooling, la network, el storage, la CPU y el accelerator, y al final solo una parte se convierte en useful work que cumple los objetivos de calidad y latencia. Esta sección convierte la discusión del ponente sobre gigawatts y facility cost en un marco de contabilidad de recursos, sin repetir cifras absolutas no auditadas.

\subsection{PUE, utilización y workload efficiency}

El \term{PUE} (Power Usage Effectiveness) se define como
\[
\mathrm{PUE}=\frac{P_{facility}}{P_{IT}},
\]
donde $P_{facility}$ es la potencia total del centro de datos y $P_{IT}$ es la IT load de los servidores, la network, el storage, etc. Cuanto más se acerca el PUE a 1, menor es el facility overhead; pero no indica si el accelerator se aprovecha plenamente ni si el modelo cumple el SLO al menor costo.

\lecturefigure{05-power-budget.png}{La facility power pasa por el overhead y la IT load antes de producir useful accelerator work.}{Subtítulos locales 10:30--14:50; esquema conceptual redibujado.}

\begin{knowledgebox}{Lectura de la figura: los tres indicadores de eficiencia son indispensables}
La \textbf{facility efficiency} se mide con el PUE; la \textbf{hardware utilization}, con el busy time del accelerator y los stalls de memoria y de network; la \textbf{workload efficiency}, con los tokens, samples o training progress obtenidos en relación con los recursos. Un clúster ocioso con un PUE bajo sigue siendo un desperdicio, y un sistema con alta utilización que genera resultados fallidos tampoco tiene valor.
\end{knowledgebox}

\begin{importantbox}{Useful compute per grid watt}
Puede usarse una expresión diagnóstica
\[
\eta_{grid}=\frac{U\cdot \eta_{workload}}{\mathrm{PUE}},
\]
donde $U$ es la tasa de utilización efectiva de los recursos de IT y $\eta_{workload}$ es la cantidad de trabajo válido producido por unidad de IT power. Aumentar $U$ puede depender del procesamiento por lotes, el scheduling y la demand aggregation; aumentar $\eta_{workload}$ puede depender de mejores kernels, model routing, quantization y cache.
\end{importantbox}

\subsection{Relojes de planificación: la electricidad, los edificios, los chips y los modelos no van sincronizados}

Lo más peligroso en los proyectos de escala nacional no es «pronosticar mal», sino amarrar entre sí decisiones de vida útil distinta. La grid y el site pueden planificarse a diez años; el facility shell requiere varios años de construcción; los accelerators y la rack density cambian con rapidez, y la model demand depende a su vez de los algoritmos y de los productos. El lector debe separar en capas la obra civil irreversible y los compute modules reemplazables.

\lecturefigure{06-planning-clocks.png}{La infraestructura nacional de IA opera al mismo tiempo en cuatro escalas temporales que entran en conflicto.}{Subtítulos locales 14:45--16:10; esquema conceptual redibujado.}

\begin{knowledgebox}{Lectura de la figura: gestionar el regret con modularity}
La capa de electricidad y terreno conserva headroom; la facility admite cooling de mayor densidad y retrofits de network; los racks y los accelerators se mantienen lo más modulares posible; el software conserva la portabilidad entre modelos y hardware. La modularity eleva el costo inicial, pero reduce la dependencia de trayectoria del tipo «cuando el centro de datos entra en operación ya no es adecuado para la nueva generación de racks».
\end{knowledgebox}

\begin{warningbox}{La capacidad planificada no es capacidad en línea}
Los 1.5--2 GW, los plazos de construcción y los costos unitarios mencionados en la sesión son cifras dadas en clase. Una evaluación formal debe distinguir entre announced, contracted, under construction, energized, installed, available y utilized capacity; de lo contrario, una misma cifra mezcla etapas distintas: el terreno, la electricidad, el building shell y los tokens que efectivamente pueden servirse.
\end{warningbox}

\teachervoice{La intuición de ingeniería del ponente es «llevar la electricidad hasta el accelerator en la mayor medida posible». Una versión más completa de los apuntes debería añadir: el accelerator solo es un activo productivo cuando lo usa una workload confiable; la capacity planning debe construir al mismo tiempo el scheduler, el software, el talento, el onboarding de modelos y el demand pipeline.}

\subsection{Resumen de la sección}

La economía de los centros de datos debe contabilizarse desde el grid watt hasta el useful outcome. El PUE, la utilization, la workload efficiency y los planning clocks determinan en conjunto el costo total; presentar cualquiera de estos indicadores por separado puede inducir a error.
